%!Tex Program = xelatex
%\documentclass[a4paper]{article}
\documentclass[a4paper]{ctexart}
\usepackage{xltxtra}
\usepackage{url}
\usepackage{booktabs}
%\setmainfont[Mapping=tex-text]{AR PL UMing CN:style=Light}
%\setmainfont[Mapping=tex-text]{AR PL UKai CN:style=Book}
%\setmainfont[Mapping=tex-text]{WenQuanYi Zen Hei:style=Regular}
%\setmainfont[Mapping=tex-text]{WenQuanYi Zen Hei Sharp:style=Regular}
%\setmainfont[Mapping=tex-text]{AR PL KaitiM GB:style=Regular} 
%\setmainfont[Mapping=tex-text]{AR PL SungtiL GB:style=Regular} 
%\setmainfont[Mapping=tex-text]{WenQuanYi Zen Hei Mono:style=Regular} 


\title{第二讲: Linux 环境基础}
\author{王何宇}
\date{}
\begin{document}
\maketitle
\pagestyle{empty}

\section{Linux 环境}

Linux 环境与其说是一些软件的使用, 更重要的是一些工作的原则和习惯. 这才是我们真正需要掌握的.
上一讲中, 我们已经体验了:

\begin{itemize}
\item 尽量使用键盘, 充分利用快捷键, 在字符界面下高效率地完成工作；
\item 用 \LaTeX 来组织你的科技文献, 使它看上去更加专业.  
\end{itemize}

不少同学因为初次接触, 因此会遇到很多实际的困难, 并且耗费了相当大的精力和时间.
但如果你要达到至少我这样一位老人家的工作效率, 这些付出是值得的. 但是千万不要忘记我们的原则:

\begin{itemize}
\item 一切都是经验和规范, 无关你的天赋和智商, 要多看文档, 多问, 善用 AI 引擎；
\item 人不要去做机器该做的事情, 不要太勤劳.
\end{itemize}

今天, 我们进一步深入这个过程. 我们先来阅读一本靠谱的书籍的部分章节:

N. Matthew and R. Stones, Beginning Linux Programming,
4th Edition, Wiley Publishing, Inc.,
Indianapolis, 2008\cite{matthew2008beginning}.



\subsection{终端和 Shell}

\textbf{终端（Terminal）} 是用户与计算机交互的一种设备或界面，用于输入命令和查看输出结果。

\begin{itemize}
    \item \textbf{历史背景：} 在早期的计算机系统中，终端是指物理设备（如键盘和显示器），
    通过串行连接到主机上。随着技术的发展，现代的 ``终端'' 通常是指运行在图形界面中的模拟程序
    （称为终端仿真器，例如 GNOME Terminal、xterm 或 Konsole）。
    \item \textbf{功能：} 提供一个窗口，允许用户输入命令并与操作系统交互；
    显示来自操作系统的输出结果。
    \item \textbf{常见终端仿真器：} GNOME Terminal、Konsole、xterm、Terminator。
\end{itemize}

\textbf{Shell} 是一个程序，充当用户与操作系统内核之间的接口。它接收用户的命令并将其传递给操作系统执行。

\begin{itemize}
    \item \textbf{作用：} 解析用户输入的命令；执行命令或将命令传递给相应的程序；
    提供脚本编程功能，允许用户自动化任务。
    \item \textbf{类型：} Shell 的种类很多，以下是一些常见的 Shell 及其特点：
\end{itemize}

\begin{table}[h]
    \centering
    \begin{tabular}{l p{0.6\textwidth}}
        \toprule
        \textbf{Shell 名称} & \textbf{特点} \\
        \midrule
        Bash & 最常用的 Linux Shell，功能强大且易于使用，支持丰富的命令和脚本功能。 \\
        Zsh & Bash 的增强版本，提供更多高级功能（如自动补全、主题支持）。 \\
        C Shell (csh) & 语法类似于 C 语言，适合熟悉 C 的开发者。 \\
        Korn Shell (ksh) & 结合了 Bourne Shell 和 C Shell 的优点，常用于商业 Unix 系统。 \\
        Fish & 用户友好的 Shell，强调易用性和现代化功能。 \\
        \bottomrule
    \end{tabular}
\end{table}

\textbf{默认 Shell：} 在大多数 Linux 系统中，默认的 Shell 是 \textbf{Bash}（Bourne Again Shell）。
可以通过命令 \texttt{echo \$SHELL} 查看当前用户的默认 Shell。

尽管终端和 Shell 经常一起使用，但它们是不同的概念：

\begin{table}[h]
    \centering
    \begin{tabular}{l p{0.4\textwidth} p{0.4\textwidth}}
        \toprule
        \textbf{特性} & \textbf{终端} & \textbf{Shell} \\
        \midrule
        定义 & 是用户与计算机交互的界面或窗口。 & 是解析和执行命令的程序。 \\
        功能 & 提供输入输出环境，将用户输入传递给 Shell。 & 接收命令并执行，返回结果给终端显示。 \\
        例子 & GNOME Terminal、Konsole、xterm。 & Bash、Zsh、C Shell、Korn Shell。 \\
        关系 & 终端是 Shell 的前端界面，Shell 是终端背后的解释器。 & \\
        \bottomrule
    \end{tabular}
\end{table}

假设你打开了一个终端窗口（如 GNOME Terminal），然后输入以下命令：

\begin{verbatim}
ls -l
\end{verbatim}

在此过程中：
\begin{itemize}
    \item \textbf{终端的作用：} 提供一个窗口让你输入命令，并显示命令的输出结果。
    \item \textbf{Shell 的作用：} 接收你输入的 \texttt{ls -l} 命令；
    解析命令并调用相应的程序（如 \texttt{/bin/ls}）；将程序的输出结果返回给终端显示。
\end{itemize}

\begin{itemize}
    \item \textbf{终端} 是用户与计算机交互的界面，提供输入输出环境。
    \item \textbf{Shell} 是一个程序，负责解析和执行用户输入的命令。
    \item 它们共同工作，使得用户可以通过字符界面的命令行与操作系统进行高效的交互。
\end{itemize}

如果一切正常，当你装完 wsl 2 之后，运行你的 ubuntu 或者 debian，你会看到一个运行着 bash 的终端.

\subsection{Linux 的基本结构}

Linux 同样具有一个文件树系统, 也就是它的文件也全部在 \verb|/| 下展开. 我们可以用 \verb|cd| 命令进入这个目录,
\begin{verbatim}
cd /
ls
\end{verbatim}

这里 \verb|ls| 是列出当前目录下全部文件的意思. 可以看到, 基本上有这样一些目录:

\begin{itemize}
\item \verb|bin|, 系统级可执行文件, 比如 \verb|ls| 这个命令就在这里；
\item \verb|usr|, 用户级系统文件, 比如用户软件. 该目录下有二级目录, 比如 \verb|/usr/bin|,
  很多应用软件的启动程序就放在这里.
\item \verb|etc|, 系统的软硬件配置文件很多都在这里.
\item \verb|home|, 用户目录, 每一个二级目录都是一个用户的全部数据. 
\item ...
\end{itemize}

理论上说, 一个用户应该只修改和维护他个人目录中的文件, {\bf 除非知道你正在干什么,
  否则不要去修改除用户目录之外的任何文件!} 即便你知道你在干什么, 切记做好适当的备份或回滚策略.
此外, 这些目录功能的分配, 只是一种约定, 有很多例外. 从现在起, 你应该尽量将数据只存放在你的个人用户目录下,
它的全称是 \verb|\home\你的用户名|. 它有一个缩写: \verb|~|. 所以, \verb|cd ~|
就可以直接进入你的个人目录. 从现在起, 要谨慎使用 \verb|sudo| 命令, 它使你作为为管理员来发布命令,
管理员可以做任何事情, 比如删除根目录: \verb|rm / -fr| ({\bf 千万不要去尝试!})

你既是你自己机器的管理员, 也是你自己机器的用户. 请时刻对你当前的角色保持清醒. 当你只想做用户时,
不要做管理员该做的事；当你想做管理员时, 确保你自己知道自己在干啥! 举个例子: 安装和配置应用软件,
并让它成为系统功能的一部分, 是典型的管理员才做的事情. 你应该一次装好, 以后尽量少折腾. 我们不培养管理员.
另外, 请千万记住, 在 Linux 下, 不要动不动就重装整个系统. 一般再严重的错误,
都可以通过部分软件的调整和安装来修复. 不要轻易用关机的方式重启系统, 这有可能导致 Linux 系统崩溃.
实际上, 如果是安装在服务器或者工作站上的 Linux 系统, 经常会很久不关机, 不重启. 甚至有计算机,
一生都没关过机. 

Linux 的命令, 实际上是一个个可执行文件. 这些可执行文件分两类, 一类是二进制可执行文件,
另一类是批处理命令, 又称 shell 脚本文件. 它是一系列 Linux 命令结合在一起完成一个复杂的功能. 
Linux 的命令之前都必须跟一个称为路径的参数. 比如真正的 \verb|ls| 的完整命令是
\verb|/bin/ls|. 但是因为系统在环境变量 \verb|PATH| 中规定了 \verb|/bin| 是一个默认的路径,
所以 \verb|/bin| 下的所有命令, 都不需要输路径前缀了. 你可以用 \verb|env | grep PATH|
来观察你当前的 \verb|PATH| 变量的内容. 比如我现在的输出是:

\begin{verbatim}
PATH=/home/hywang/anaconda3/bin:/home/hywang/anaconda3/condabin
:/usr/local/sbin:/usr/local/bin:/usr/sbin:/usr/bin:/sbin:/bin
:/usr/games:/usr/local/games:/snap/bin
\end{verbatim}

这些目录之下的可执行文件, 都不需要路径前缀. 但是注意, 当前目录 \verb|.| 并不在路径中. 因此,
即便你就在命令所在的目录下, 但是这个目录不在路径中, 你也必须打路径前缀. 这就是我们学习 C 语言时,
得到的可执行文件, 运行必须加 \verb|./运行文件| 的原因.


退回上一层目录的命令是 \verb|cd ..|, 这里 \verb|..| 是上一层目录的意思.
你可以一直退到 \verb|/|.

Linux 的二进制可执行文件是怎么来的? 它其实就是 C 语言编译出来的. 当你构建文件 \verb|hello.c|:

\begin{verbatim}
#include <stdio.h>
#include <stdlib.h>
int main()
{
    printf(“Hello World\n”);
    exit(0);
}
\end{verbatim}

然后运行:

\begin{verbatim}
gcc -o hello hello.c
\end{verbatim}

以后, 你就得到一个可执行文件 \verb|hello|, 它本质上和一个 Linux 系统命令没有任何区别. 
你如果够无聊, 可以将它复制到缺省路径 \verb|/usr/local/bin| 下,  

\begin{verbatim}
sudo cp hello /usr/local/bin
\end{verbatim}

然后你就可以在任何目录下运行一个新命令 \verb|hello|. 

最后, 介绍一个很有用的命令: \verb|info|. 它们提供 Linux 命令的官方文档. 比如:
\verb|info ls|. 但现在，你只有在没有网络的时候才需要用它。

\section{用 shell 脚本来加速}

首先介绍一个 UNIX 哲学: 一次做好一件具体的事情, 然后将零星的功能链接起来, 成为一个更加丰富的命令.
或者说: 作为系统架构师, 你在创建具体命令时, 要考虑到更高层次上的重用；而作为用户,
尽可能利用现有的命令的链接去完成更复杂的功能. 用互联网黑话说, 就是不要重复造轮子.

\subsection{管道和重定向}
比如: \verb|ls -l| 是列出当前目录下的全部文件, 但这样会导致一屏幕可能放不下.
\verb|more| 命令则是将一个超出一屏幕的输出, 在输出一满屏内容后暂停, 然后可以逐行(按回车)输出.
将两个命令用 \verb!|! 管道(pipe) 链接在一起: \verb!ls -l | more! 就是将当前目录下所有文件列出,
按屏停顿.

我们之前还有一个例子, \verb|env| 实际上是输出全部环境变量, 而 \verb|grep| 是过滤的意思,
所以 \verb!env | grep PATH! 就是输出全部环境变量中包含 \verb|PATH| 内容的东西.

这些组合都可以灵活搭配使用. UNIX 哲学在所有分支中都是通行的. 

大家可能注意到, shell 本身也是一个命令. 事实上, 目前流行的 shell 有多个版本, 我们使用的是 bash,
这也是开源社区最通行的版本. 使用 Mac OS 的同学可能注意到 Mac 的默认 shell 不是 bash.
这个命令本身的位置在 \verb|/bin/bash|. Linux 的命令实际上在 shell 中运行和反馈, shell 则提供 pipe
和环境变量等等这样的运行环境. 除了 pipe, shell 还有一个重要特性是重定向(redirection).

我们知道, C 语言中输出, 用 \verb|printf| 函数, 本质上是向一个被称为 \verb|stdout|
的设备文件输出字符流. \verb|printf| 和 \verb|stdout| 都在头文件 \verb|stdio.h| 
定义. 在正常情况下, \verb|stdout| 是定义在显示器上的. 但我们可以用重定向, 将其重定向为一个文本文件.
比如:
\begin{verbatim}
ls -l > re.txt
\end{verbatim}
将当前目录的内容重定向到 \verb|re.txt| 下. 这里 \verb|>| 每一次都会重新生成一个新的输出结果文件.
而 \verb|>>| 则是在文件结尾新增.
\begin{verbatim}
ls -l >> re.txt
\end{verbatim}

有的时候, 我们写了一个程序, 会有大量的输出, 如果从屏幕走, 记录本身就很麻烦.
这时你没有必要专门做一套文本文件的读写机制, 而只需要重定向就行了. 除了 stdout 可以重定向,
stderr 也是可以重定向的. 为了区别, 我们可以分别用 1(可以忽略) 和 2 表示这两个不同的通道. 

\subsection{shell 编程}

我们可以做的更过份一点. 比如我们知道在 \verb|/usr/include| 下放了很多 C 语言的头文件,
比如 stdio.h 就在这里. 现在我们想看一下这些头文件中有多少包含了 \verb|stdout| 这个关键字.
比如我们想看一下这个文件是在哪个文件定义的. 那么我们可以用命令:

\begin{verbatim}
grep stdout /usr/include/*
\end{verbatim}

找到这些文件, 然后记下来, 再一个个文件去读. 这就有点蠢了. 就不能一次搞定? 找到这些包含
\verb|stdout| 的头文件, 找到一个, 读一个, 这样就很快能发现是谁定义, 还能马上看到定义周围的细节.
要做到这件事, 我们可以写一个段小的脚本:

\begin{verbatim}
$ for file in /usr/include/*
> do
> if grep -l stdio $file
> then
> more $file
> fi
> done
\end{verbatim}

这样就可以愉快地做到了. 注意我们可以用上箭头调用上一次运行的命令, 
用 TAB 键自动补齐一个最有可能的命令.


\section{例子讲解}

\begin{verbatim}
#!/bin/sh

foo=”Hello World”
echo $foo
unset foo
echo $foo
\end{verbatim}

\verb|unset| 的功能是解除了变量 \verb|foo| 的设置.

\section{正则表达式的应用}

正则表达式(regular expression) 是定义字符串模式（pattern）的一种规则.
凡是满足指定规则的字符串, 我们就说和该正则表达式匹配(match). 
大家其实早就接触过一些正则表达式规则, 比如通配符 \verb|*| 和 \verb|?|.
前者表示不限长度的任何字符, 后者表示任何单个字符.

详细的正则表达式规则, 请参见: \newline
https://www.runoob.com/regexp/regexp-tutorial.html

我们直接看一些 shell 脚本中的实际例子. 

\section{find 和 grep}

我们经常遇到的问题是不知道一个文件在哪里. 比如我们知道修改 apt 的源需要修改 source.list 文件,
但有时因为 linux 版本什么的原因, 我们不知道这个文件到底在哪里. 这时可以用
\begin{verbatim}
sudo find / -name sources.list -print
\end{verbatim}
从根目录开始以超级用户身份搜索可以确保你不会漏掉任何地方,
但如果你插了 U 盘, 也许你不想搜索进入 U 盘(因为你很确信你要找的是一个系统文件), 于是可以用参数
\verb|-mount| 避免搜索全部挂载的文件系统:
\begin{verbatim}
sudo find / -mount -name sources.list -print
\end{verbatim}
对比运行, 你可以发现现在快了很多, 并且避开了一些奇奇怪怪的地方.
完整的 \verb|find| 命令是这样的:
\begin{verbatim}
find [path] [options] [tests] [actions]
\end{verbatim}
这里 \verb|path| 是开始查找的路径；\verb|options| 就是比如 \verb|-mount|
这样的选项, \verb|tests| 实际上比正则表达式更加丰富一些, 而 \verb|actions|
可以让你指定查找之后的工作.

这里注意第 65 页例子中 \verb|{}| 和 \verb|\;| 的作用\cite{zhoudx2020universality}.

而 \verb|grep| 的作用是在文间内找关键词, 当然也可以是一堆文件内. 命令格式
\begin{verbatim}
grep [options] PATTERN [FILES]
\end{verbatim}
这里的 \verb|PATTERN| 是正则表达式. \verb|options| 资料上列举了几个重要的:

\begin{table}
  \centering
  \begin{tabular}{|c|c|}
    \hline
    \verb|options| & 功能 \\
    \hline
    -c & 只输出有几行匹配\\
    \hline
    
  \end{tabular}
\end{table}

\section{项目文件和系统目录}
现在我们注意到, 一个合理规划的项目, 它下面的目录结构首先应该井然有序. 一般来说, 遵从一些默认的约定:

\begin{table}[!phb]
  \centering
  \begin{tabular}{|c|l|}
    \hline
    \bf{目录名} & \bf{存放内容} \\
    \hline
    \verb|doc| & 文档, 说明. \\
    \hline
    \verb|src| & 源代码, \verb|cpp| 后缀. \\
    \hline
    \verb|include| & 头文件, \verb|h| 后缀. \\
    \hline
    \verb|lib| & 库文件, \verb|a|, \verb|o|, \verb|so| 后缀. \\
    \hline
    \verb|bin| & 可执行文件. \\
    \hline
    \verb|example| & 例子. \\
    \hline
  \end{tabular}
  \caption{常见的目录约定}
  \label{tab::dir}
\end{table}

之所以会出现这种结构分布, 和 Unix 体系下 C\/C++ 的规则有关. 我们知道,
C\/C++ 的风格是主程序短小精悍, 而实际功能由各个函数完成. 函数部分有分为声明,
也就是定义和数据结构和函数的使用形式和接口；而函数的具体实现则又放在具体的库文件中.
库文件是一种预先编译好的二进制代码. 它的形式一般是 \verb|.o|, \verb|.a| 或者
\verb|.so| 后缀的文件. 这里 \verb|.o| 就是一个源码的直接编译后版本. 比如我们可以尝试一下:
\begin{verbatim}
g++ -o bitmap.cpp
\end{verbatim}
我们就看到对应生成了一个 \verb|bitmap.o| 文件. 将这个文件和 \verb|bitmap.h| 文件一起交给用户,
用户就可以在他的程序中使用 \verb|bitmap.h| 中定义的各种函数了, 此时用户并不需要获得
\verb|bitmap.cpp| 文件. 这样有助保持这个程序的稳定性.
因为避免了用户直接去改源码所导致的不可知问题. 有时我们会将多个功能接近的函数的二进制文件打包成一个库文件,
这样用户得到一个库文件就可以使用其中的全部函数. 这种文件, 如果后缀是 \verb|a|, 则称为静态库；
如果后缀是 \verb|so|, 称为动态库. 它们的区别是, 当用户调用静态库时,
这部分函数二进制代码就一起嵌入了用户的最终可执行文件；而用户调用动态库时,
大部分函数二进制代码不会嵌入最终可执行文件真正在运行时, 只有在这个可执行文件真正被执行时, 由操作系统临时对接,
完成函数功能. 动态库的好处是生成的最终代码比较小, 而且不容易发生代码泄露, 相对更加安全.
因此我们现在基本上都采用动态库. 尝试一下:

\begin{verbatim}
g++ -c -fPIC bitmap.cpp
g++ -shared -o libbitmap.so bitmap.o
\end{verbatim}

这里 \verb|-fPIC| 的含义是生成一个相对寻址的二进制文件. 这个和汇编结构有关, 这里不详述.
因为是动态库, 所以要能够在内存中根据相对地址而不是绝对地址运行. 上面的命令第一行产生
\verb|bitmap.o| 文件. 而第二行将此文件转换成动态链接库, 也即 \verb|libbitmap.so| 文件.
动态链接库总是以 \verb|lib*.so| 的形式命名, 也是一种约定.

现在我们可以将 \verb|bitmap.h| 和 \verb|libbitmap.so| 分别放到 \verb|~\include|
和 \verb|~\lib| 的目录下(只要做链接就可以, 如果没有目录可以直接建立),
这样我们就可以在个人目录的范围下使用 \verb|bitmap| 了.
\begin{verbatim}
ln -s ~/Projects/MathSoft/src/bmp/bitmap.h ~/include
ln -s ~/Projects/MathSoft/src/bmp/libbitmap.so ~/lib
\end{verbatim}

我们可以在 \verb|~\tmp| 下尝试一下看看.
\begin{verbatim}
g++ -o test test_bmp.cpp -I../include -L../lib -lbitmap
\end{verbatim}
这里 \verb|-I| 给出了头文件位置, 而 \verb|-L| 给出了库文件的位置.
\verb|-l| 则指明要调用哪个库, 这里的规则是 \verb|-l| 加上库文件
\verb|libxxx.so| 中的 \verb|xxx| 部分, 也即 \verb|-lxxx|.
然后我们顺利地得到了可执行文件 \verb|test|, 但当我们运行时, 会得到错误:
\begin{verbatim}
./test: error while loading shared libraries: libbitmap.so: 
cannot open shared object file: No such file or directory
\end{verbatim}

这个错误充分体现了动态库的调用机制, 在生成可执行文件时, 只是标记一下,
并不检查文件是否正常. 只有在运行时, 它才会去寻找这个 \verb|so| 文件并试图执行.
但是从效率出发, 它并不是直接去读文件系统, 而是从系统维护的一个数据库中去检索.
因为我们还没有注册我们的 \verb|libbitmap.so| 文件, 所以它找不到.

现在来注册, 我们在 \verb|/etc/ld.so.conf.d/| 这个目录下新建一个文本文件
\verb|hywang.conf| (你可以用你自己的名字), 内容只有一行: 
\begin{verbatim}
/home/hywang/lib
\end{verbatim}
这样以后凡是存放或者链接到这个目录的 \verb|.so| 文件, 都会被系统自动找到.
系统更新数据库会有一个时间, 我们手动更新一下:
\begin{verbatim}
sudo ldconfig
sudo ldconfig -p | grep bitmap
\end{verbatim}
第一行是更新数据库, 第二行则是检查一下数据库里是否已经有我们的 \verb|libbitmap.so|
文件了. 现在回到 \verb|~/tmp| 下, 不用重新编译, 我们就发现我们 \verb|test|
文件可以运行了. 

所以一个系统级别的应用, 将会在系统的头文件目录和库文件目录留下头文件和动态链接库.
比如我们熟悉的 \verb|math.h| 文件, 当我们调用它时, 需要 \verb|#include <math.h>|,
然后在编译时, 需要 \verb|-lm|. 现在我们知道, 系统中除了 \verb|#include <math.h>| 以外,
必然还存在 \verb|libm.so| 文件. 用我们学过的 \verb|find| 技巧将它们找出来.
\begin{verbatim}
find / -mount -name math.h -print
find / -mount -name libm.so -print
\end{verbatim}

现在我们总结一下, 一个项目, 在开发完毕后, 项目目录下 \verb|include| 里存放头文件,
\verb|src| 里存放源文件(\verb|cpp| 文件). 然后编译时, 在\verb|lib| 产生 \verb|so| 文件,
在 \verb|bin| 产生可执行文件. 最后在安装阶段, 将全部头文件复制(或链接)到用户的系统头文件目录
(比如 \verb|\usr\local\include|), 将库文件复制(或链接)到用户的系统库文件目录
(比如 \verb|\usr\local\lib|), 然后用户就可以在全系统愉快地使用了.

以上工作, 颇为繁琐, 且容易出错. 这就需要引出我们的一个项目自动生成工具: \verb|make|.

\section{自动编译脚本 make}

就这个 bitmap 的项目, 我们现在来做个简单的整理. 我们需要做的事情是:

\begin{enumerate}
\item 在 \verb|lib| 目录中生成 \verb|libbitmap.so| 文件,
  要用到的文件是 \verb|bitmap.h|;
\item 在 \verb|bin| 目录中生成 \verb|test| 文件,
  要用到的文件是 \verb|bitmap.h|, \verb|libbitmap.so|;
\item 在用户指定的目录中创建头文件和库文件的链接,
  用户的可执行目录中创建测试的可执行文件. 
\end{enumerate}

这里的每一步都涉及若干行满是参数的命令, 为了避免手工出错,
Unix\/Linux 社区提供了多个解决方案, 最常见的之一是 \verb|make| 命令.
如果没有安装 \verb|make|, 请用 \verb|apt| 安装:
\begin{verbatim}
sudo apt install make
\end{verbatim}

命令 \verb|make| 首先要提供一个特殊的脚本文件, 一般命名为 \verb|Makefile|.
这个文件中指定了编译的规则. 基本规则是:
\begin{verbatim}
目标 : 所需文件
    脚本命令
\end{verbatim}
各目标之间如果有依赖, 则会根据依赖自动回溯, 确保依赖完整. 比如目标 A 需要目标 B 所生成的文件,
那么当目标 B 的内容发生变化时, \verb|make| 会确保目标 A 也被重新编译.

现在来写一个最简单的 \verb|Makefile|. 我们先将 \verb|bitmap| 项目的目录结构调整为
\begin{verbatim}
include/bitmap.h
src/bitmap.cpp
\end{verbatim}

我们进入 \verb|bitmap| 项目目录, 构建 \verb|Makefile| 文件, 内容为:
\begin{verbatim}
libbitmap.so : bitmap.o
	g++ -shared -o libbitmap.so bitmap.o

bitmap.o : src/bitmap.cpp include/bitmap.h
	g++ -c -fPIC src/bitmap.cpp -Iinclude 
\end{verbatim}
然后在 \verb|bitmap| 目录下执行 \verb|make|, 我们看到, 自动编译开始生成 \verb|bitmap.o|
和 \verb|libbitmap.so| 文件. 如果我们删除 \verb|libbitmap.so|, 再打一次 \verb|make|,
我们看到只会重新生成 \verb|libbitmap.so| 文件. 但是如果删除 \verb|bitmap.o| 文件, 再打一次
\verb|make|, 这时发现两个文件都会重新生成. 因为前者依赖后者, 当后者变化时, 前者也不可靠,
必须重新生成以保证依赖. 当你有上千个文件互相依赖时, 这一机制就非常重要了.

文件 \verb|Makefile| 中的脚本命令本质上就是 Shell 脚本, 所以我们实际上可以干很多事. 比如,
我们希望创建的库文件都放在 \verb|lib| 目录下, 而可执行文件都放在 \verb|bin| 目录下,
我们可以这样修改 \verb|Makefile|.
\begin{verbatim}
lib/libbitmap.so : src/bitmap.cpp include/bitmap.h 
	g++ -c -fPIC src/bitmap.cpp -Iinclude 
	g++ -shared -o libbitmap.so bitmap.o
	@if [ ! -d lib ]; then mkdir lib; fi
	@mv libbitmap.so lib
	rm bitmap.o
\end{verbatim}

这里我们不再把中间文件 \verb|bitmap.o| 作为一个目标, 而且每次编译完成,
还将这个文件删除了, 保持目录干净. 注意我们这里用 Shell 脚本的 \verb|if|
命令判定是否已经建立了 \verb|lib| 目录.

我们还可以建立没有依赖的功能性目标, 比如
\begin{verbatim}
clean:
	rm lib/libbitmap.so
\end{verbatim}
这样我们只要打 \verb|make clean| 就可以恢复编译前的样子.
但这里有一件尴尬的事情, 如果正好有一个文件名字叫 \verb|clean|, 比如我们执行
\begin{verbatim}
touch clean
\end{verbatim}
现在 \verb|make clean| 就会认为是需要更新这个文件, 得不到正确的响应. 解决的办法是做个标记:
\begin{verbatim}
.PHONY: clean
\end{verbatim}
这种标记称为伪目标, 就是不管怎样, 都必须执行的目标. 现在我们完善这个编译的构建:
\begin{verbatim}
bin/test : src/test_bmp.cpp include/bitmap.h lib/libbitmap.so
	g++ -o test src/test_bmp.cpp -Iinclude -Llib -lbitmap 
	@if [ ! -d bin ]; then mkdir bin; fi
	@mv test bin

lib/libbitmap.so : src/bitmap.cpp include/bitmap.h 
	g++ -c -fPIC src/bitmap.cpp -Iinclude 
	g++ -shared -o libbitmap.so bitmap.o
	@if [ ! -d lib ]; then mkdir lib; fi
	@mv libbitmap.so lib
	rm bitmap.o

clean:
	rm lib/libbitmap.so bin/test bin/*.bmp

install: lib/libbitmap.so bin/test
	@if [ ! -d ~/include ]; then mkdir ~/include; fi
	@if [ ! -d ~/lib ]; then mkdir ~/lib; fi
	@if [ ! -d ~/bin ]; then mkdir ~/bin; fi
	ln -s lib/libbitmap.so ~/lib
	ln -s include/bitmap.h ~/include
	ln -s bin/test ~/bin

.PHONY: clean install
\end{verbatim}

\section{注释和文档}

代码必须有规范的文档. 否则这段代码将不可维护, 也很难共享. 程序必须有文档, 否则用户将无法使用.
在 Linux 社区中, 这两件事情往往通过一个被称为 doxygen 的工具统一完成.
首先请确认安装了 doxygen:
\begin{verbatim}
sudo apt install doxygen
\end{verbatim}
然后如果有可能的话, 开启你的编辑软件的辅助工具. Emacs 下就是 doxmacs, 而 vscode 则可以在它的
Market Place 找到各种支持. doxymacs 已经停止维护, 并且不在 apt 源里, 所以不得不手动安装.

doxygen 巧妙利用注释符号的自由空间, 以特殊格式的注释符号对程序注释进行标记. 因为程序注释和文档具有共性,
因此有可能将两件事情在一个框架内解决.

doxygen 的文档可以从 apt 获得:
\begin{verbatim}
sudo apt install doxygen-doc
\end{verbatim}

doxygen 支持 latex, 因此可以引入公式. doxygen 可以中文化, 会有一些困难, 但不是不可以克服.



\bibliographystyle{plain}
\bibliography{crazyfish.bib}

\end{document}
